% Fichier complété par tools/complete_missing_corrections.py.
% Source : GEO/GEO-03/18_Maxpid/18_Maxpid.tex
% Statut : rédigé (5 question(s)).
\begin{corrigebox}[Corrigé rédigé]
\CorrigeQuestion{1}
Le graphe des liaisons est
                $0\xleftrightarrow[\text{pivot}]{D}4\xleftrightarrow[\text{pivot}]{C}3$,
                avec $3\xleftrightarrow[\text{glissière}]{\vec x_1}1$,
                $1\xleftrightarrow[\text{pivot}]{B}0$ et la liaison hélicoïdale
                $2/3$ de pas $p=4\,\mathrm{mm}$. Le rotor 2 est guidé en pivot dans le
                stator 1. La chaîne géométrique utile est donc le triangle fixe-mobile
                $BCD$, dont $BC=\lambda$ et $DC=r$ (longueur portée par le levier 4).
\CorrigeQuestion{2}
En notant $L=BD$, $r=DC$ et $\theta_0$ l'angle fixe de
                $\overrightarrow{BD}$ avec $\vec x_0$, la fermeture du triangle $BCD$
                donne
                $$\lambda^2=L^2+r^2-2Lr\cos(\theta-\theta_0).$$
                Sur la branche correspondant à la configuration dessinée,
                $$\boxed{\theta(\lambda)=\theta_0+
                \arccos\!\left(\frac{L^2+r^2-\lambda^2}{2Lr}\right)}.$$
                Les constantes $L$, $r$ et $\theta_0$ se calculent une fois pour toutes
                avec les dimensions $a,b,c,d$ du schéma.
\CorrigeQuestion{3}
La dérivation de la fermeture fournit
                $2\lambda\dot\lambda=2Lr\sin(\theta-\theta_0)\dot\theta$.
                Ainsi, hors position singulière,
                $$\boxed{\dot\theta=
                \frac{\lambda}{Lr\sin(\theta-\theta_0)}\,\dot\lambda}.$$
                Le signe est automatiquement celui de la branche géométrique choisie.
\CorrigeQuestion{4}
Pour une vis de pas $p$, un tour du rotor produit une translation $p$ :
                $$\dot\lambda=\frac{p}{2\pi}\,\omega.$$
                Par conséquent
                $$\boxed{\dot\theta=
                \frac{p\lambda}{2\pi Lr\sin(\theta-\theta_0)}\,\omega}.$$
                Cette expression met en évidence l'amplification de vitesse à proximité
                des alignements, qui sont des configurations singulières.
\CorrigeQuestion{5}
À $n=500\,\mathrm{tr\,min^{-1}}$, la vitesse moteur vaut
                $\omega=2\pi n/60=52{,}36\,\mathrm{rad\,s^{-1}}$ et
                $\dot\lambda=pn/60=33{,}3\,\mathrm{mm\,s^{-1}}$ pour $p=4\,\mathrm{mm}$.
                Un tracé Python robuste est obtenu en balayant les valeurs admissibles de
                $\lambda$, puis en calculant $\theta$ et $\dot\theta$ :
                \begin{verbatim}
import numpy as np
import matplotlib.pyplot as plt
p = 4e-3
n = 500/60
omega = 2*np.pi*n
# L, r, theta0 : valeurs déduites des cotes a, b, c, d
lam = np.linspace(abs(L-r)+1e-5, L+r-1e-5, 1000)
alpha = np.arccos((L**2+r**2-lam**2)/(2*L*r))
theta = theta0 + alpha
theta_dot = (p/(2*np.pi))*omega*lam/(L*r*np.sin(alpha))
plt.plot(theta, theta_dot)
plt.xlabel(r'$\theta$ (rad)'); plt.ylabel(r'$\dot\theta$ (rad/s)')
plt.grid(); plt.show()
                \end{verbatim}
                On écarte volontairement les extrémités, où $\sin(\theta-\theta_0)=0$.
\end{corrigebox}
